\vssub
\subsubsection{~Wave height assimilation with optimal interpolation ({\code DA1})} \label{sub:da1}
\opthead{DA1}{WAM model, UWA\footnotemark[12], BOM\footnotemark[13]}{M. Derkani, S. Zieger}

\footnotetext[12]{University of Western Australia}
\footnotetext[13]{Australian Bureau of Meteorology}
\noindent
Data assimilation package ({\F w3wda1} in {\file w3wda1md.ftn}) features
an optimal interpolation scheme for assimilating wave height data
into {\ws}. The scheme is based on optimal interpolation for
assimilating altimeter data from \citet{art:LGJ92},
originally implemented in WAM \citep{bk:WAM94}.

The analysis wave height field is estimated by optimim interpolation using
the correlation lengh scale calculation of either \citet{art:VMH97}, or
\citet{GY2004}. {\ws}'s full two-dimensional wave spectrum is utilised as
first-guess spectrum and classified into primarily wind sea (or swell)
using the models internal partitioning algorithm (see topographic partitioning
in section \ref{sub:num_part}) to determine transformation parameters $A$ and $B$.
Knowing the first-guess spectrum $F_1(f,\theta)$, the analysis spectrum
$F_\mathrm{a}(f,\theta)$ is constructed with Eq.~(\ref{eq:DA101}).:
\begin{equation}\label{eq:DA101}
F_\mathrm{a}(f,\theta) = A\cdot F_1(B\cdot f,\theta)
\end{equation}
Depending on wind sea or swell spectrum, different techniques are applied
to estimate tranformation parameters from analysis (subscript `a`) and
first-guess spectrum (subscript `1`). 

If the first-guess spectrum contains primarily wind sea, parameters $A$ and $B$ 
can be determined in accordance with nondimensional growth laws, specifically
the nondimensional energy $E\star=E\,g^2/u_\star^4$, the nondimensional
mean frequency $\overline{f}\star = u_\star\,\overline{f}/g$ and the nondimensional
duration $t\star=g\,t/u_\star$ relation. Knowing parameters (i.e., $H_s$, $u_\star$, and $\overline(f)$)
of the wind sea partition from the first-guess spectrum, an analysis friction
velocity and analysis mean frequency are determined from duration limited
growth relation \citep{art:LGJ92} (Eq.~(\ref{eq:DA104})) and 
frequency-energy growth \citep{art:TGLDAS22} curve (Eq.~(\ref{eq:DA105})).
\begin{eqnarray}
E\star = 955 \tanh {\bigl ( 6.02\cdot 10^{-5}\,t\star^{0.695}\bigr )}\label{eq:DA104}\\
f\star = 5.054\cdot 10^{-4}\,\overline{f}\star^{-2.959}\label{eq:DA105}
\end{eqnarray}

Parameters $A$ and $B$ are then calculated from significant wave height $H_s$ 
and mean frequency $\overline{f}$ as given in Eq.~(\ref{eq:DA102}).
\begin{equation}\label{eq:DA102}
A = \left ( \frac{H_{s\,\mathrm{a}}}{H_{s\,1}} \right )^2\text{,}\ 
B = \frac{\overline{f}_\mathrm{a}}{\overline{f}_1}
\end{equation}

In case the spectrum contains primarily swell, parameters
$A$ and $B$ are calculated from Eq.~(\ref{eq:DA103}) with a small correction $C$ applied.
\begin{equation}\label{eq:DA103}
C = 1.0 - 6\cdot 10^{-3} \bigl ( H_{s\,\mathrm{a}} - H_{s\,1} \bigr )\text{,}\ 
B = C \left (  \frac{H_{s\,\mathrm{a}}}{H_{s\,1}} \right )^{1/2}\text{,}\ 
A = B \left ( \frac{H_{s\,\mathrm{a}}}{H_{s\,1}} \right )^{2}
\end{equation}

% tab:DA1_par

\begin{table} \begin{center}
\begin{tabular}{|l|c|l|} \hline \hline
  Namelist  &  value   & description \\
\hline
   METHOD   &  0    &  \citet{art:VMH97} (default) \\
            &  1    &  \citet{GY2004} \\
   MAXDKM   & 4000  &  distance (in km) around observation to activate routine \\
   SEAFCUT  &  0.5  &  wind sea dominance threshold \\
  \hline \hline
\end{tabular} \end{center}
\caption{Namelist parameters for data assimilation package {\code \&WDA1}.}
\label{tab:DA1_par} \botline \end{table}


